\begin{textsample}{POS Dim 5 – human – Score -1.00 – sp/sp\_inf\_e\_ef\_intro\_7\_p11.txt}  \label{ex:f5_pos_007}
Na Geografia, por exemplo, é comum o uso do termo alfabetização cartográfica para referir -se a um conjunto de saberes e de fazeres relacionados a noções básicas, como o reconhecimento de área e sua representação, identificação da visão vertical e oblíqua presentes em mapas, da linha, do ponto, da escala da proporção, a leitura de legendas, o reconhecimento de imagens bidimensionais e tridimensionais, a orientação e a utilização e leitura dos pontos de referências, entre outros, fundamentais para desenvolver a autonomia na leitura e na produção de representações do espaço.

% matched lemmas: 
\end{textsample}
